

The corresponding metrics are detailed in \autoref{tab:geographic_causal}.

\begin{table}[H]
\centering
\caption{Summary of Supplementary Quasi-Experimental Estimates}
\label{tab:geographic_causal}
\small
\begin{tabular}{@{}llrrrr@{}}
\toprule
\textbf{Design} & \textbf{Dep.\ Var.} & \textbf{Coeff.} & \textbf{SE} & \textbf{$p$} & \textbf{$N$} \\
\midrule
\textit{(1) 10-1 Council ITS} & vote\_no & \metricTenOneITSCoeff{} & 0.122 & \metricTenOneITSPval{} & 20 \\
\textit{(2) NPA Designation} & vote\_no & \metricNPAFrictionCoeff{} & 0.245 & \metricNPAFrictionPval{}& 20 \\
\textit{(3) Historic District} & vote\_no & \metricHDFrictionCoeff{} & --- & \metricHDFrictionPval{} & 20 \\
\textit{(4) TOD Deregulation} & vote\_no & \metricTODFrictionCoeff{} & --- & \metricTODFrictionPval{}& 20 \\
\textit{(5) 2022 Flipped District DiD}& protested & \textbf{\metricFlipDiDCoeff{}} & 0.266 & \textbf{\metricFlipDiDPval{}} & 518 \\
\bottomrule
\multicolumn{6}{l}{\footnotesize OLS with robust standard errors. Designs (3)--(4) are non-informative due to limited variation}\\
\multicolumn{6}{l}{\footnotesize or zero-variance outcomes; they are reported for transparency rather than substantive interpretation.} \\
\multicolumn{6}{l}{\footnotesize Design (5) reports the $\text{Flipped} \times \text{Post-2022}$ interaction coefficient.} \\
\end{tabular}
\end{table}